% =====================================================================
% Capítulo 2 — Esboço de seção/interlúdio (análise de rede textual)
% Modelo: esboco_secao_infranodus.tex (capítulo 1). O método já foi
% apresentado no interlúdio do capítulo 1; aqui retomo o gesto e leio a
% rede específica deste capítulo.
%
% Estilo seguido (regras da tese): primeira pessoa do singular, sem
% travessões em texto principal, "capítulo" minúsculo em referências
% cruzadas, aspas com \enquote{}, sem fórmulas do tipo "não X, mas Y".
%
% Observação técnica: para o build encontrar a figura, copie
%   infranodus/cap2/infranodus_cap2_focus.png
% para
%   figuras/cap.2/infranodus_cap2_focus.png
% (mantendo a convenção do capítulo). As entradas bibliográficas usadas
% aqui já foram sugeridas no rodapé de esboco_secao_infranodus.tex.
% =====================================================================

\subsection*{Interlúdio: a rede textual do capítulo}

Retomo neste ponto o gesto de análise de rede textual que apresentei no
capítulo~\ref{capitulo1}, aplicando à íntegra deste capítulo o mesmo
\textit{pipeline} de co-ocorrência, comunidades de Louvain
\parencite{blondel2008fast} e centralidade
\parencite{paranyushkin2019infranodus}. Convertidos os
28.174~\textit{tokens} do capítulo em um grafo de 180 termos, a
detecção de comunidades separa três vocabulários que reconheço como as
três operações do argumento. O primeiro é a leitura cerrada das
figurações em Latour, e o grafo o registra de um modo que me interessa:
as próprias obras aparecem condensadas em pares lexicais de altíssima
associação, \textit{laboratory}~$\leftrightarrow$~\textit{life}
(NPMI~0,87), \textit{hope}~$\leftrightarrow$~\textit{pandora}
(0,87) e \textit{action}~$\leftrightarrow$~\textit{science}
(0,79), de modo que \enquote{Laboratory Life}, \enquote{A Esperança
de Pandora} e \enquote{Science in Action} entram na rede como
etiquetas lexicais compactas. O segundo vocabulário é a genealogia
militar-industrial da inteligência artificial, ancorada em
\textit{militar}, \textit{indústria}, \textit{inteligência} e
\textit{artificial}. O terceiro é o panorama bibliométrico do campo
brasileiro, reunido em torno de \textit{capes}, \textit{distribuição},
\textit{produção} e \textit{corpus}.

Os termos de maior \textit{degree} ponderado são \textit{artificial},
\textit{inteligência}, \textit{latour}, \textit{militar} e
\textit{ciência}, e o par \textit{inteligência}~$\leftrightarrow$~%
\textit{artificial} é a co-ocorrência mais pesada de todo o capítulo
(peso~819, NPMI~0,83): o objeto que descrevo aparece nomeado
sempre da mesma forma, colado, quase uma palavra só. A maior ponte do
grafo, por \textit{betweenness}, é \textit{latour} (0,32),
seguido de \textit{inteligência}, \textit{artificial}, \textit{ciência},
\textit{militar} e \textit{vocabulário}. Que \textit{latour} seja o termo
que mais costura comunidades distintas é coerente com um capítulo que
faz do vocabulário latouriano o seu operador de leitura; e que
\textit{vocabulário} apareça ele próprio como ponte devolve ao capítulo
a sua própria imagem, já que o gesto do capítulo é um exercício sobre
vocabulários, sobre o que as figurações fazem quando nomeiam.

\begin{figure}[!htb]
\centering
\includegraphics[width=1.0\textwidth,keepaspectratio]{figuras/cap.2/infranodus_cap2_focus.png}
\caption[Rede textual do capítulo~\ref{capitulo2} (\textit{text network
analysis})]{Rede textual do capítulo~\ref{capitulo2}
(\textit{text network analysis}). Os nós são os termos de maior
\textit{degree} ponderado; o tamanho do nó acompanha o \textit{degree};
a cor indica a comunidade detectada por Louvain ponderado; as arestas,
arqueadas, recebem a mistura das cores das comunidades que ligam.
Versão de alta resolução, arquivo \texttt{.gexf} para o \textit{Gephi} e
relatório completo em \texttt{infranodus/cap2/} no repositório desta
tese.}
\label{fig:infranodus-cap2}
\end{figure}

O que o grafo me devolve com mais força é uma assimetria que sustenta o
argumento do capítulo. O polo militar-industrial está lexicalmente denso:
\textit{militar} figura entre os termos de maior \textit{degree} e entre
as maiores pontes, e articula uma vizinhança própria. A contrafiguração
têxtil-feminista que oponho a esse polo, ao contrário, mal se inscreve na
superfície da rede, e sua ausência entre os termos centrais é ela mesma
um dado: a figuração militar já está sedimentada no vocabulário com que o
campo fala de si, enquanto a alternativa haraweana ainda precisa ser
tecida no texto para deixar rastro mensurável. Leio essa desproporção
como confirmação empírica modesta da tese que o capítulo defende, e como
tarefa, já que tornar a contrafiguração audível é torná-la também
lexicalmente presente.

A segunda leitura vem das ligações mais fracas. As comunidades que menos
se costuram são a do panorama bibliométrico (\textit{capes},
\textit{distribuição}, \textit{produção}) e a das obras de Latour lidas
no original (\textit{science}, \textit{action}, \textit{pandora},
\textit{hope}); fraca também é a costura entre o vocabulário da
inteligência artificial e essas mesmas obras. O mapeamento quantitativo
do campo e a leitura cerrada das figurações correm em paralelo sem se
tocarem, e a rede torna visível uma junta que o próprio capítulo se
propõe a fazer trabalhar: mostrar como a lexicometria das obras-fonte
conversa com o panorama bibliométrico da literatura que as recebe. Os
actantes menores que o grafo isola, o par
\textit{crawford}~$\leftrightarrow$~\textit{joler} (0,74) da
anatomia da IA e a etiqueta \textit{aime} do inquérito latouriano sobre
os modos de existência, confirmam que a rede registra com fidelidade as
referências que mobilizo, ainda quando as deixo em posição periférica.
